(1) 非相对论作用量 \(S=\int(mv^2/2-mgz)\,dt\)。两水平臂速率相同，动能相位抵消，重力势差为 \(mgH\)，传播时间 \(L/v\)，故相位差大小 \(|\Delta\phi|=mgHL/(\hbar v)=mgA/(\hbar v)\)；符号取决于相位差定义。
(2) 由 \(v=h/(m\lambda)=2\pi\hbar/(m\lambda)\)，得到 \(|\Delta\phi|=m^2gA\lambda/(2\pi\hbar^2)\)。这是重力势能差累积的量子作用量相位，不是中子在路径上实际获得额外重力势能。
(3) 若面积法向与重力方向的夹角使有效竖直投影面积为 \(A\sin\alpha\)，则 \(\Delta\phi(\alpha)=mgA\sin\alpha/(\hbar v)\)。当倾角变化满足 \(\Delta(\sin\alpha)=2\pi\hbar v/(mgA)\) 时条纹移动一个周期；若右端大于可达投影范围则不能完成一整周期。
(4) \(g\to0\)、\(A\to0\) 或 \(v\to\infty\) 时重力相位趋零。推导假定两臂的水平速率相同、连接段的动力学相位对称抵消，并忽略晶体散射相位差；真实实验需校准这些系统效应。 还应分离地球自转的Sagnac相位，以免误判重力效应。
